\documentclass{beamer}
\usepackage[german]{babel}
%\usepackage{libertine}
%\renewcommand*\familydefault{\sfdefault}  %%
\usepackage[default]{opensans}
%\usepackage{gillius2}
%\usepackage{lato}
\usepackage{siunitx}

\usepackage[utf8]{inputenc}

\usepackage[T1]{fontenc}

\usepackage{booktabs} % To thicken table lines
\usepackage{multirow}


\usepackage{microtype}
\usepackage{xcolor,multido,colortbl}
\usepackage{pgffor}

\usepackage{tikz}
\usepackage[version=4]{mhchem}
\usepackage{tikzorbital}

\usepackage{pgfplots}
 %\pgfplotsset{com  =1.13}
\usepackage{graphicx}
\usepackage{adjustbox}
\usepackage{hyperref}
\usepgfmodule{matrix} 
\usetikzlibrary{matrix,positioning}
\usetikzlibrary{arrows}

\newcommand{\makemycolor}[2]{%
    \pgfmathsetmacro{\hue}{(#1/100)^1.715*0.79}%
    \definecolor{myhsbcolor}{hsb}{\hue,1,1}%
    \textcolor{myhsbcolor}{#2}%
}
\newcommand\gauss[2]{1/(#2*sqrt(2*pi))*exp(-((x-#1)^2)/(2*#2^2))} % Gauss function, parameters mu and sigma

\usepackage[absolute,overlay]{textpos}
%\usecolortheme[light,accent=red]{solarized}
%\definecolorseries{H}{hsb}{last}[hsb]{0,1,1}[hsb]{1,1,1}\resetcolorseries[10]{H}
%\definecolorseries{S}{hsb}{last}[hsb]{0.1,0,0}[hsb]{.1,1,1}\resetcolorseries[10]{S}
%\definecolorseries{B}{hsb}{last}[hsb]{1,1,0}[hsb]{1,1,1}\resetcolorseries[10]{B}
\tikzset{
    level1/.style = {
        ultra thick,
        green
    },
    level4/.style = {
        ultra thick,
        blue
    },
    level5/.style = {
        ultra thick,
        red
    }
}

\newcommand\ytl[2]{
\parbox[b]{8em}{\hfill{\color{cyan}\bfseries\sffamily #1}~$\cdots\cdots$~}\makebox[0pt][c]{$\bullet$}\vrule\quad \parbox[c]{4.5cm}{\vspace{7pt}\color{red!40!black!80}\raggedright\sffamily #2.\\[7pt]}\\[-3pt]}
%\usepackage[texcoord,grid,gridunit=mm,gridcolor=red!10,subgridcolor=green!10]{eso-pic}
\setbeamertemplate{navigation symbols}{}


\title{Lumineszenz von Seltenerd-Ionen}
\author{Nevroz Arslan}

\begin{document}
  {%
    \setbeamertemplate{headline}{}
   \setbeamercolor{postit}{fg=black,bg=yellow}
    \frame{\titlepage}
  }

\begin{frame}[t]\frametitle{Was macht ein Leuchtstoff gut?}
Es kommt nicht darauf an!
\begin{itemize}
   \item hohe Quantenausbeute $QE = \frac{\text{emmitierte \ Photonen}}{\text{absorbierte Photonen}}\%$
   \item hohe Übergangswahrscheinlichkeit, kurze Lebensdauer und scharfe Emmision 
   \item geringe Neigung zu strahlungsloser Relaxation \ce{Ln^* -> Ln +} Wärme
\end{itemize}
   
\end{frame}

\begin{frame}[t]\frametitle{Was macht ein Leuchtstoff bunt?}
\begin{columns}
   \column{.7\textwidth}
    \begin{itemize}
    \item Farbe ist nicht gleich Wellenlänge
    \item Pink im sichtbaren Spektrum? 
  % \item Lila ist dunkles Purpur.
    \end{itemize}
   \column{.3\textwidth}
      \begin{tikzpicture}
    \foreach \k in {0,1,...,100}{%
        \pgfmathsetmacro{\hue}{(\k/100)^1.715*0.79}
        \definecolor{mycolor}{rgb:hsb}{\hue,1,1}
        \node[color=mycolor] () at (0,\k/20) {$\bullet$};
    }%
    \foreach \f in {0,1,...,10}{%
        \pgfmathtruncatemacro{\num}{(\f*30)+480}
        \node () at (0.75,\f/2) {$\num$ nm};
    }
    \foreach \g/\h in {0/Rot,2/Orange,4/Gelb,6/Grün,8/Blau,10/Purpur}{%
        \pgfmathtruncatemacro{\num}{\g*10}
        \node at (-1,\g/2) {\makemycolor{\num}{\h}};
    }%
    \end{tikzpicture}

\end{columns}

\end{frame}


 
%Ein Farbraum beschreibt die Menge der darstellbaren Farben eines Mediums 

\begin{frame}[t]\frametitle{Die drei natürlichen Koordinaten der Farbe}
    

\begin{itemize}
      \item Farbton (Wellenlänge) %\multido{\nColr=0+1}{10}{\fcolorbox{white}{H!![\nColr]}{}}
      \begin{tikzpicture}
          \colorlet{color min hsb}[hsb]{red}
  \colorlet{color max hsb}[hsb]{magenta}
   \foreach \pos in {100,...,0}{
    \colorlet{my color hsb}[rgb]{color min hsb!\pos!color max hsb}
    \fill[fill=my color hsb,draw=none] (-3+\pos/20,0) rectangle +(5mm,3mm);
  }
      \end{tikzpicture}
   \item Sättigung (Weißanteil) 
\begin{tikzpicture}
          \colorlet{color min hsb}[hsb]{white}
  \colorlet{color max hsb}[hsb]{red}
   \foreach \pos in {100,...,0}{
    \colorlet{my color hsb}[rgb]{color max hsb!\pos!}
    \fill[fill=my color hsb,draw=none] (-3+\pos/20,0) rectangle +(5mm,3mm);
  }
      \end{tikzpicture}
      \item Helligkeit (Intensität ) %\multido{\nColr=0+1}{10}{\fcolorbox{white}{B!![\nColr]}{}}
      \item Strahlungsleistung (Wiens oder Planck) 
      \item $\lambda _{max} \cdot T $ = konst.
   \end{itemize}
 \end{frame}



%Mit zunehmender Temperatur verschiebt sich die maximale Strahlungsintensität eines Schwarzen Körpers zu kürzeren Wellenlängen, der Farbeindruck wechselt dabei vom Roten ins Bläulich-Weiße. Der Farbton einer (Wärme-) Lichtquelle lässt sich als Temperatur eines vergleichbaren Schwarzen Strahlers angeben. Damit erhält man die Farbtemperatur der Lichtquelle.
\begin{frame}[t]\frametitle{Farbetemperatur}
\begin{beamerboxesrounded}[shadow=false]{}
\begin{figure}[!h]
\centering
      \includegraphics[width=\textwidth]{pics/ct.png}
      \caption{\footnotesize Farbtemperatur nach dem planckschen Strahlungsgesetz}
 \end{figure}
\end{beamerboxesrounded}

\begin{figure}[!h]
\centering
\begin{tikzpicture}[scale=0.8]

\begin{axis}[every axis plot post/.append style={
  mark=none,samples=50,smooth}, % All plots: from -2:2, 50 samples, smooth, no marks
axis x line*=bottom, % no box around the plot, only x and y axis
axis y line*=left, % the * suppresses the arrow tips
enlargelimits=upper] % extend the axes a bit to the right and top
%\addplot table[x=x,y=y] {pics/probe.dat};
\end{axis}
\end{tikzpicture}
 \end{figure}
\end{frame}
   
\begin{frame}[t]\frametitle{Was macht ein Leuchtstoff umweltfreundlich?}
  \begin{beamerboxesrounded}[]{Umweltfreunlich}
  \begin{itemize}
   \item geringere UV-Anteile im Licht 
  \end{itemize}
  \end{beamerboxesrounded}
 
\end{frame}
 %Unter den Insekten sind vor allem Nachtfalter betroffen, deren Augen stark auf die UV-Strahlung von Hochdrucklampen reagieren. Die Gebäude- und Straßenlampen werden stundenlang umflogen, bis die Tiere ermatten oder vor Entkräftung sterben.https://de.wikipedia.org/wiki/Quecksilberdampflampe

\begin{frame}[t]\frametitle{Die Farben der Lanthanoiden}

\begin{beamerboxesrounded}[shadow=false]{}
  Die Lanthanoide treten bevorzugt dreiwertig auf.
\end{beamerboxesrounded}
     \begin{table}
      \centering
\resizebox{\linewidth}{!}{%
\begin{tabular}{cccccccccccccccc}
\toprule
 $[Xe]$ &\ce{La^{3+}} & \ce{Ce^{3+}} &\cellcolor[RGB]{ 110,194,47} \ce{Pr^{3+}} &\cellcolor[RGB]{202,99 ,138} \ce{Nd^{3+}} & \cellcolor[RGB]{ 240,42 ,242} \ce{Pm^{3+}}&\cellcolor[RGB]{  252, 255, 110} \ce{Sm^{3+}} & \ce{Eu^{3+}} & \ce{Gd^{3+}} & \ce{Tb^{3+}}& \cellcolor[RGB]{ 163, 255, 0} \ce{Dy^{3+}} &  \cellcolor[RGB]{  255,255, 0} \ce{Ho^{3+}} & \cellcolor[RGB]{  194,171, 187}\ce{Er^{3+}} & \cellcolor[RGB]{171,228,149} \ce{Tm^{3+}} & \ce{Yb^{3+}} & \ce{Lu^{3+}}\\
   & \cellcolor[RGB]{254,225,2} \ce{Ce^{4+}} & \cellcolor{yellow}\ce{Pr^{4+}} &\cellcolor[RGB]{ 111, 0,255} \ce{Nd^{4+}} &&  &  & \cellcolor[RGB]{151,26,14} \ce{Sm^{2+}} & \ce{Eu^{2+}} & \cellcolor[RGB]{ 251, 241, 35}\ce{Dy^{4+}}&  &  & &  &\cellcolor[RGB]{ 163,79,128}\ce{Tm^{2+}}&  \cellcolor[RGB]{181,201,19}\ce{Yb^{2+}}\\
   &  & &  &  &  &  &  &  & \cellcolor[RGB]{60,0,0} \ce{Tb^{4+}} &  &  &  &  & & \\
\midrule
  4f & 0 &1 & 2 & 3 & 4 & 5 & 6 & 7 & 8 & 9 & 10 & 11 & 12 & 13& 14\\
  \bottomrule
\end{tabular}
} 
\end{table}
 
    \begin{itemize}
      \item \ce{La} (\ce{f^0}), \ce{Gd} (\ce{f^7}), \ce{Lu} (\ce{f^{14}}) sind nur dreiwertig.
      \item Die stabilen f-Konfiguration werden erreicht in \ce{Ce^{4+}}, \ce{Tb^{4+}}, \ce{Eu^{2+}},\ce{Yb^{4+}}
    \end{itemize}
\end{frame}


\begin{frame}[t]\frametitle{Schalenaufbau}
    
   \begin{columns}
    \column{.8\textwidth}
      
    \begin{adjustbox}{max totalsize={.9\textwidth}{\textheight},center}
\begin{tikzpicture}
\draw [->,ultra thick] (0,0) --   (0,10) node[above] {\Large Energie};
\draw [->,ultra thick] (0,0) -- node[below] {\Large n (Schale)} (20,0);
    \draw[level4] (6,1) -- node[above] {4s} (7,1);
       \draw[level4] (6,2.5) -- node[above] {4p} (7,2.5);
       \draw[level4] (4.8,2.5) --  (5.8,2.5);
       \draw[level4] (7.2,2.5) --  (8.2,2.5);
      \draw[level4] (6,4) -- node[above] {4d} (7,4);
       \draw[level4] (7.2,4) --  (8.2,4);
       \draw[level4] (8.4,4) --  (9.4,4);
       \draw[level4] (4.8,4) --  (5.8,4);
       \draw[level4] (3.6,4) --  (4.6,4);
       \draw[level4] (6,6) -- node[above] {4f} (7,6);
       \draw[level4] (7.2,6) --  (8.2,6);
       \draw[level4] (8.4,6) --  (9.4,6);
       \draw[level4] (4.8,6) --  (5.8,6);
        \draw[level4] (3.6,6) --  (4.6,6);
        \draw[level4] (2.4,6) --  (3.4,6);
      \draw[level4] (9.6,6) --  (10.6,6);
      \draw[level5] (14,3)  -- node[above] {5s} (15,3);
       \draw[level5] (14,5) -- node[above] {5p} (15,5);
       \draw[level5] (12.8,5) --  (13.8,5);
       \draw[level5] (15.2,5) --  (16.2,5);
       \draw[level5] (14,6.5) -- node[above] {5d} (15,6.5);
       \draw[level5] (12.8,6.5) --  (13.8,6.5);
       \draw[level5] (15.2,6.5) --  (16.2,6.5);
        \draw[level5] (11.6,6.5) --  (12.6,6.5);
       \draw[level5] (16.4,6.5) --  (17.4,6.5);
       \draw[level1] (18,5.9) -- node[above] {6s} (19,5.9);
    \end{tikzpicture}
\end{adjustbox}
    \column{.2\textwidth}
     
\end{columns}

\end{frame}


\begin{frame}[t]\frametitle{RusselSanderskopplung}
\begin{itemize}
  \item Mehrelektronenzustände
  \item für einen gegebenen Wert von L existieren jeweils $2L+1$ entarteter Bandzustände 
  \item charakterisiert durch die magnetischen $M_L = 3,2,1,0,-1,-2,-3$ 
  \item S = Spinquantenzahl 
  \pause
  \item \begin{tikzpicture}[scale=0.8]

\drawLevel[elec = up, pos = {(0,3)},    width = 1]{ndf1};
\drawLevel[elec = up, pos = {(1.2,3)},  width = 1]{ndf2};
\drawLevel[elec = up, pos = {(2.4,3)},  width = 1]{ndf3};
\drawLevel[pos = {(3.6,3)},  width = 1]{ndf4};
\drawLevel[pos = {(4.8,3)},  width = 1]{ndf5};
\drawLevel[pos = {(6,3)},  width = 1]{ndf6};
\drawLevel[pos = {(7.2,3)},  width = 1]{ndf7};
\drawLevel[pos = {(8.4,3)},  width = 1]{ndf8};
\node[right] at (right ndf8) {\Large \ce{^{4}_{}I_{9/2} Nd^{3+}  } } ;

\end{tikzpicture}
\end{itemize}
%\ce{Y2O3:Eu^{2+}} bedeutet, dass ein Wirtgitter aus Yttriumoxid mit Europium-Ionen dotiert ist
% d. h. dass ca. fünf Prozent der Yttrium-Ionen durch Europium-Ionen ersetzt wurde.

\end{frame}

\begin{frame}[t]\frametitle{Interkombinationsverbot\cite{holeman}}
    %jeder Übergang, bei dem sich der sich Gesamtspin geändert, ist verboten 
\begin{figure}[!h]
\centering
\begin{tikzpicture}[scale=0.8]

\drawLevel[elec = updown, pos = {(0,3)},    width = 1]{ndf1};
\drawLevel[elec = up, pos = {(1.2,3)},  width = 1]{ndf2};
\drawLevel[elec = up, pos = {(2.4,3)},  width = 1]{ndf3};
\drawLevel[elec = up, pos = {(3.6,3)},  width = 1]{ndf4};
\drawLevel[elec = up, pos = {(4.8,3)},  width = 1]{ndf5};
\drawLevel[elec = up, pos = {(6,3)},  width = 1]{ndf6};
\drawLevel[elec = none, pos = {(7.2,3)},  width = 1]{ndf7};
\node[right] at (right ndf7) {\Large \ce{^{5}_{}S_{7/2} Gd^{3+}  } } ;


\drawLevel[elec = up, pos = {(0,1)},    width = 1]{ndf21};
\drawLevel[elec = up, pos = {(1.2,1)},  width = 1]{ndf2};
\drawLevel[elec = up, pos = {(2.4,1)},  width = 1]{ndf3};
\drawLevel[elec = up,pos = {(3.6,1)},  width = 1]{ndf4};
\drawLevel[elec = up,pos = {(4.8,1)},  width = 1]{ndf5};
\drawLevel[elec = up,pos = {(6,1)},  width = 1]{ndf6};
\drawLevel[elec = up,pos = {(7.2,1)},  width = 1]{ndf9};
\node[right] at (right ndf9) {\Large \ce{^{8}_{}S_{7/2} Gd^{3+}  } } ;
\node at (0.5,4) {$-3$};
\node at (1.5,4) {$-2$};
\node at (2.7,4) {$-1$};
\node at (4,4) {$0$};
\node at (5.5,4) {$1$};
\node at (6.5,4) {$2$};
\node at (7.5,4) {$3$};
\draw [<-, ultra thick] (left ndf1) edge[bend right]  (left ndf21);
\end{tikzpicture}
\caption{Interkombinationsverbot}

\end{figure}
\begin{itemize}
  \item Verboten, weil zu Spinpaarungen führt.
  \item die farblosen Beispiele aus ÜM: \ce{Fe^{3+}}, \ce{Mn^{2+}}
  \item farblos auch \ce{La^{3+} f^0}, \ce{Gd^{3+} f^7},\ce{Lu^{3+} f^{14}}   
\end{itemize}
\end{frame}
\begin{frame}[t]\frametitle{Regel von Laporte}
\begin{itemize}
  \item Übergänge zwischen unterschiedlicher Parität erlaubt
  \item s- und d-Orbitale sind gerade
  \item p- und d-Orbitale sind ungerade  
\end{itemize}
  
\end{frame}

\begin{frame}[t]\frametitle{\small Derzeit eingesetzte Leuchtstoffe für Color-80-Lampen.}
    \begin{tabular}{lccccc}
\toprule
\multirow{2}{*}{Leuchtstoff}& Abk. &  \ce{$\lambda$} & Abs. bei  & QE & LE  \\
& &\si{\nano\meter}&254 \si{\nano\meter}  & \% & \si{\lumen\per\watt}\\
\midrule
 \ce{BaMgAl_{10}O_{17}:Eu}& BAM & \cellcolor[wave]{450} 450  & 90 &  90 &90\\
\ce{(Ce,Tb)MgAl_{11}O_{19}}   & CAT   &  \cellcolor[wave]{541} 541  &  95 & 90& 495\\
\ce{Y_2O_3:Eu}    & YOX   &  \cellcolor[wave]{611} 611  &  75 & 90 &280\\
\ce{LaPO_4:Ce,Tb}  & LAP  & \cellcolor[wave]{545} 545 & 95 &93 & 500\\
\bottomrule
\end{tabular}

\end{frame}

\begin{frame}[t]\frametitle{Konzentration Abhängigkeit }
\begin{itemize}
  \item Mit Leuchtstoffen gezielt verunreinigte Halbleiterkristalle
  \item It has been found that the average life period is practically independent of temperature and wavelength but depends markedly on concentration
  \item the intensity, polarisation and the average life period are all dependent on concentration
  \item 
\end{itemize}


\end{frame}

  \begin{frame}
  \frametitle{Quantenteiler}
 
  \end{frame}

\begin{frame}
  \frametitle{Stokes Verschiebung}
    \begin{itemize}
    \item 
    \end{itemize}
  \end{frame}


\begin{frame}[t]\frametitle{Was macht ein Leuchtstoff gut?}
\begin{itemize}
 \item Anregungsenergie möglichst effizient absorbieren
 \item Quantenausbeute möglichst hoch sein soll
\end{itemize}
\end{frame}



\begin{frame}[t]\frametitle{Literaturverzeichnis}
    

\begin{thebibliography}{9}

\bibitem{wikilila}
\emph{Violett}
 \url{https://de.wikipedia.org/wiki/Violett}.
\bibitem{holeman}
E. Wiberg, \emph{Lehrbuch der Anorganische Chemie}, 102. Aufl., De Gruyter, Berlin \textbf{2007}, S. 1370.
\end{thebibliography}
  
\end{frame}


%http://www.rsc.org/chemistryworld/News/2009/September/01090901.asp
%https://www.osapublishing.org/oe/fulltext.cfm?uri=oe-19-S3-A331
\end{document}

% Reference